\documentclass[11pt]{article}
\usepackage{series}
\usepackage{mathtools}

\title{Admissibility Precosheaves and Conditional AQFT Nets\\
in Modal Triplet Theory}
\author{Peter Nero}
\date{July 2026}

\hypersetup{
  pdftitle={Admissibility Precosheaves and Conditional AQFT Nets in Modal Triplet Theory},
  pdfauthor={Peter Nero}
}

\begin{document}
\maketitle

\begin{abstract}
This paper replaces the former claim that algebraic quantum field theory
(AQFT) follows from admissible charts alone.  Two structures must be kept
separate.  An admissibility-indexed precosheaf is a pregeometric covariant
functor only after its algebras and extension morphisms are supplied; overlap
or failure of joint representability does not create commutators or imply
locality.  A physical Haag--Kastler net instead requires a selected Lorentzian
base, an upper local operator net, and a localization-preserving coherent
reduction.  For a decomposable orthogonal projector $P$, we prove that
compression of the $P$-compatible upper subalgebras preserves isotony and
spacelike commutation.  This is the rigorous MTT-to-AQFT bridge currently
available.  We also show why absence of a global admissible chart does not
forbid an abstract quasilocal algebra or categorical colimit.  It may obstruct
a single chart-induced physical realization, but states, faithful
representations, horizon behavior, and irreversibility require separate
theorems.  The result is a conditional locality-descent theorem and a precise
ledger of the data still required for a full AQFT reconstruction.
\end{abstract}

\section*{Revision note for this edition}
\addcontentsline{toc}{section}{Revision note for this edition}
\begin{description}
\item[Supersedes.] Version~1, \emph{From Modal Triplet Theory to Algebraic
Quantum Field Theory: Local Nets from Admissible Charts and Coherent Basin
Persistence}.
\item[Reason.] Version~1 inferred commutation from non-joint
representability, assumed extension of every smaller-chart observable, and
identified absence of a global chart with absence of an abstract global
algebra.
\item[Resolution.] Version~2 distinguishes the pregeometric chart functor from
the physical Haag--Kastler net and bases physical locality on an upper local
net plus coherent locality descent.
\item[Retained result.] Admissible charts can index a useful partial algebraic
description when their algebra objects and transition morphisms are explicitly
given.
\item[Open boundary.] A selected upper QFT net, state space, covariance,
spectrum/positivity, time-slice property, nonperturbative continuum limit, and
the chart-to-region intertwiner remain independent obligations.
\end{description}

\tableofcontents

\section{Scope and correction of the chart-only argument}

AQFT assigns algebras to causally localized spacetime regions and relates them
by injective algebra morphisms \cite{HaagKastler1964}.  Locally covariant QFT
packages analogous data as a covariant functor on globally hyperbolic
spacetimes and causal embeddings \cite{BFV2003}.  Neither framework obtains
Einstein causality merely from the absence of a common coordinate chart.

The former edition tried to derive three conclusions directly from MTT
admissible charts:
\begin{enumerate}
\item inclusion of charts was said to extend every observable and hence prove
isotony;
\item failure of joint representability was said to force commutation; and
\item absence of one global chart was said to forbid a global algebra.
\end{enumerate}
All three steps are invalid without additional structure.  Extension is a
map that must be defined and proved compatible.  A commutator is meaningful
only after both operands are represented in a common algebra.  An abstract
inductive limit or colimit need not be represented by any one indexing
object.

The corrected construction therefore has three stages:
\[
\begin{aligned}
 \text{admissibility precosheaf}
 &\longrightarrow \text{selected Lorentzian upper net}\\
 &\longrightarrow \text{coherently compressed physical net}.
\end{aligned}
\]
The first arrow requires an explicit chart-to-region and algebra
intertwiner.  The second is the locality-descent theorem proved below.

\section{The AQFT target}

Let $(Y,g)$ be a time-oriented globally hyperbolic Lorentzian spacetime, and
let $\mathcal K(Y)$ be a chosen poset of relatively compact causally convex
open regions.  Write $\mathbf{Alg}$ for a category of unital $*$-algebras (or
$C^*$-algebras) and unital injective $*$-homomorphisms.

\begin{definition}[Physical local net]
A local net is a covariant functor
\[
 \mathfrak A:\mathcal K(Y)\longrightarrow\mathbf{Alg}.
\]
For an inclusion $O_1\subseteq O_2$, functoriality supplies an injective map
$\iota_{12}:\mathfrak A(O_1)\to\mathfrak A(O_2)$.  Identifying its image with
a subalgebra gives isotony,
\[
 O_1\subseteq O_2
 \quad\Longrightarrow\quad
 \mathfrak A(O_1)\subseteq\mathfrak A(O_2).
\]
\end{definition}

\begin{definition}[Einstein causality]
The net is local if, whenever $O_1$ and $O_2$ are spacelike separated,
\[
 [\mathfrak A(O_1),\mathfrak A(O_2)]=0
\]
inside a declared common ambient algebra.  For a graded field net, the ordinary
commutator is replaced by the graded commutator.
\end{definition}

These two conditions are only part of a physical AQFT.  Depending on the
target, one must also address covariance, a state space, spectrum or
microlocal-spectrum conditions, additivity, the time-slice axiom, duality,
superselection structure, and a suitable continuum dynamics.  The time-slice
property is a dynamical theorem even in established perturbative
constructions, rather than a consequence of isotony alone
\cite{ChilianFredenhagen2008}.

If the region system is directed and the morphisms are injective, its
$C^*$-inductive limit gives a quasilocal algebra
\[
 \mathfrak A_{\mathrm{ql}}
 =\varinjlim_{O\in\mathcal K(Y)}\mathfrak A(O).
\]
The limit is an abstract universal object.  It need not equal
$\mathfrak A(O_*)$ for a largest region $O_*$, and no largest region need
exist.

\section{The pregeometric admissibility category}

Let $\mathcal C_{\mathrm{adm}}$ be a category whose objects are declared MTT
admissible contexts
\[
 \alpha=(U_\alpha,P_\alpha,\mathcal D_\alpha,\mathbf m_\alpha),
\]
where $U_\alpha$ is a domain, $P_\alpha$ a coherent reduction,
$\mathcal D_\alpha$ the required operator-domain data, and
$\mathbf m_\alpha$ a vector of admissibility margins.  Raw set inclusion
$U_\alpha\subseteq U_\beta$ is not yet an algebra morphism.

\begin{definition}[Admissibility precosheaf]
An admissibility precosheaf is a covariant functor
\[
 \mathfrak B:\mathcal C_{\mathrm{adm}}\longrightarrow\mathbf{Alg}.
\]
A morphism $u:\alpha\to\beta$ includes, or is represented by, a specified
unital injective $*$-homomorphism
\[
 j_u:\mathfrak B(\alpha)\longrightarrow\mathfrak B(\beta)
\]
satisfying $j_{\operatorname{id}}=\operatorname{id}$ and
$j_{v\circ u}=j_v\circ j_u$.
\end{definition}

This definition is useful but conditional.  The chart domains and projectors
do not determine $\mathfrak B(\alpha)$, and overlap alone does not determine
$j_u$.

\begin{proposition}[No automatic extension]\label{prop:noextension}
An inclusion $U_\alpha\subseteq U_\beta$ and admissibility of both charts do
not imply that every observable in $\mathfrak B(\alpha)$ extends to
$\mathfrak B(\beta)$, nor that any extension is unique or injective.
\end{proposition}

\begin{proof}
Extension is additional algebraic data.  The same pair of domains can be
assigned a constant algebra functor, inequivalent matrix algebras, quotient
algebras, or algebras with no selected homomorphism between them.  None of
these choices is fixed by set inclusion or positivity of admissibility margins.
If a restriction map from the larger domain is supplied, it runs in the
opposite direction and still need not possess a section.
\end{proof}

\begin{proposition}[Non-comparability is not commutation]
Suppose $\alpha$ and $\beta$ have no common comparison object in
$\mathcal C_{\mathrm{adm}}$.  This fact alone implies neither
$[A,B]=0$ nor $[A,B]\ne0$ for $A\in\mathfrak B(\alpha)$ and
$B\in\mathfrak B(\beta)$.
\end{proposition}

\begin{proof}
Without morphisms into a common algebra, the product $AB$ and hence the
commutator $AB-BA$ are not typed.  A value cannot be inferred for an expression
that has not been defined.  If a later physical realization maps both
observables into a common algebra, their commutator is determined by that
realization and its locality theorem.
\end{proof}

Thus $\mathfrak B$ records which partial descriptions and extensions have
actually been constructed.  It is a pregeometric organization of
representability, not yet a Haag--Kastler net.  In particular it has no notion
of spacelike separation until a causal localization functor is supplied.

\section{Physical realization data}

The physical construction requires the following independent hypotheses.

\begin{description}
\item[Lorentzian base.] A selected, time-oriented, globally hyperbolic
four-dimensional base $(Y,g)$ and a bundle $\pi:M\to Y$.

\item[Upper local theory.] A concrete upper net
\[
 O\longmapsto\mathfrak A_U(\pi^{-1}O)\subseteq\mathcal B(\mathcal H_U)
\]
that is isotonic and local on $\mathcal K(Y)$.

\item[Coherent reduction.] A decomposable orthogonal projector
\[
 P=\int_Y^\oplus P_y\,\mathrm d\nu(y)
\]
on the upper Hilbert bundle.  Decomposability prevents the reduction itself
from mixing unrelated base fibers.

\item[Compatibility.] Physical observables are taken from the subalgebra that
preserves the coherent sector,
\[
 \mathfrak A_U^P(O)
 :=\{A\in\mathfrak A_U(\pi^{-1}O):[A,P]=0\}.
\]
For a graded net, $P$ is also required to be even.
\end{description}

FP~VI does not construct these data from the fixed-point spine alone.  It
proves that an instantaneous spatially bilocal kernel is insufficient for
causality and identifies a local hyperbolic mediator theory as the consistent
completion route \cite{FPVIv4}.  A selected completion must still supply the
upper local net.  The corrected Projection--Admissibility paper then supplies
the compression theorem \cite{Projectionv2}.

\section{Coherent locality descent}

For $O\in\mathcal K(Y)$ define
\begin{equation}
 \mathfrak A_P(O)
 :=\{PAP|_{P\mathcal H_U}:A\in\mathfrak A_U^P(O)\}
 \subseteq\mathcal B(P\mathcal H_U).
 \label{eq:compressed-net}
\end{equation}

\begin{theorem}[Conditional MTT locality descent]\label{thm:descent}
Under the physical realization hypotheses above, $O\mapsto\mathfrak A_P(O)$
is an isotonic local net.  If the upper net is $C^*$-algebraic, each
$\mathfrak A_P(O)$ is a $C^*$-algebra.  The graded statement also holds when
$P$ is even and upper locality is graded locality.
\end{theorem}

\begin{proof}
For $A,B\in\mathfrak A_U^P(O)$, commutation with $P$ gives
\[
 (PAP)(PBP)|_{P\mathcal H_U}=PABP|_{P\mathcal H_U},
 \qquad
 (PAP)^*=PA^*P.
\]
Thus compression is a unital $*$-homomorphism on the compatible subalgebra,
and its image is a unital $*$-algebra.  A $*$-homomorphic image of a
$C^*$-algebra is $C^*$-algebraic after the standard quotient/image
identification.

If $O_1\subseteq O_2$, upper isotony gives
$\mathfrak A_U^P(O_1)\subseteq\mathfrak A_U^P(O_2)$, hence
$\mathfrak A_P(O_1)\subseteq\mathfrak A_P(O_2)$.  If $O_1$ and $O_2$ are
spacelike separated and $A,B$ belong to their respective compatible upper
algebras, then
\[
 [PAP,PBP]|_{P\mathcal H_U}
 =P[A,B]P|_{P\mathcal H_U}=0
\]
by upper locality.  Replacing commutators by graded commutators proves the
graded case.
\end{proof}

The theorem is inheritance, not creation.  It cannot prove locality unless
the upper theory is already local, and it says nothing about an upper
observable that fails to preserve $P\mathcal H_U$.  Likewise, a nondecomposable
projector may mix base support and lies outside the declared physical
interpretation.

\section{The missing chart-to-region intertwiner}

The pregeometric and physical constructions are connected only after a
selected interface is provided.  Let
\[
 L:\mathcal K(Y)\longrightarrow\mathcal C_{\mathrm{adm}}
\]
assign an admissible context to each physical region.  Suppose there are
$*$-homomorphisms
\[
 \rho_O:\mathfrak B(L(O))\longrightarrow\mathfrak A_P(O).
\]

\begin{proposition}[Naturality gate]\label{prop:naturality}
The admissibility precosheaf represents the physical net only if, for every
$O_1\subseteq O_2$, the naturality equation
\begin{equation}
 \rho_{O_2}\circ j_{12}
 =\iota_{12}\circ\rho_{O_1}
 \label{eq:naturality}
\end{equation}
holds.  If every $\rho_O$ is an isomorphism, the two nets are naturally
isomorphic.  If they are merely homomorphisms, the interface is a reduction or
representation, not an equivalence.
\end{proposition}

\begin{proof}
Commutativity is exactly the naturality condition.  Componentwise
isomorphisms define a natural isomorphism; weaker components do not.
\end{proof}

Equation~\eqref{eq:naturality}, together with a selected $L$, is the precise
replacement for the former assertion that chart overlap automatically becomes
physical localization.

\section{Global algebras, states, and representations}

Absence of a terminal or largest object in $\mathcal C_{\mathrm{adm}}$ means
that no one admissible chart represents every context.  It does not imply
that the diagram $\mathfrak B$ lacks a categorical colimit, and it does not
imply that the physical net lacks a quasilocal algebra.

\begin{proposition}[Compatible states and the inductive limit]
Let $\{\mathfrak A_P(O),\iota_{12}\}$ be a directed unital $C^*$-net.  A state
$\omega$ on $\mathfrak A_{\mathrm{ql}}$ restricts to a compatible family
$\omega_O$.  Conversely, a compatible family satisfying
\[
 \omega_{O_2}\circ\iota_{12}=\omega_{O_1}
\]
defines a state on the algebraic inductive limit and hence on its
$C^*$-completion.
\end{proposition}

\begin{proof}
Restriction gives compatibility by functoriality.  Conversely, define the
functional on an equivalence class represented by $A\in\mathfrak A_P(O)$ as
$\omega_O(A)$.  Compatibility makes the value independent of the
representative.  Positivity, normalization, and boundedness pass to the
completion.
\end{proof}

Therefore a no-state claim must prove failure of compatible physical state
data, not merely failure of a global chart.  Similarly, every abstract
$C^*$-algebra has a faithful Hilbert-space representation, while a
\emph{selected chart-induced representation satisfying all MTT admissibility
conditions} may fail to exist.  Those are different statements.

\section{Horizons and irreversibility are separate targets}

A physical AQFT can assign algebras to regions on both sides of a causal
horizon.  Restriction to an observer's wedge, a thermal or modular property,
and loss of operational access depend on the spacetime, net, and state.  A
boundary of an MTT admissible chart may obstruct continuation of that chart or
of one selected representation, but it becomes a physical horizon only after
the map $L$ and the relevant causal statement are proved.

Likewise, neither a missing global chart nor a failed extension map establishes
an arrow of time.  Irreversibility needs oriented dynamics plus an asymmetric
property such as noninvertible effective evolution, a monotone entropy or
Lyapunov functional, or boundary data.  The local-net theorem is neutral on
that question.

\section{Current MTT-to-AQFT status}

The current corpus supports the following scoped ledger.

\begin{description}
\item[Admissibility-indexed precosheaf.] Available as a conditional
pregeometric construction once algebra objects and extension morphisms are
specified.  It is not forced by chart overlap alone.

\item[Physical Lorentzian base.] The selected q79 branch currently supplies a
conditional global Lorentzian coframe and causal representative after the
discrete $A_{\mathrm{QG}}$ realization declaration and one binary
$A_{\mathrm{causal}}$ boundary mark.  This does not itself construct a QFT
net.

\item[Upper local theory.] FP~VI identifies a local hyperbolic parent as the
valid causal completion route but does not select its field content, gauge
algebra, couplings, state, or operator net.

\item[Locality descent.] Theorem~\ref{thm:descent} is exact under its stated
upper-net and projector hypotheses.  This is the currently rigorous
AQFT-style bridge.

\item[Perturbative observable branch.] A current selected-SM result gives a
conditional perturbative observable functor, while importing standard
BRST/Faddeev--Popov quantization rather than deriving it from MTT.

\item[Constructive finite-domain QFT.] Current SPT-filtered TT/BRST functional
integrals provide conditional finite-domain Borel and Ward-identity results.
Infinite volume, the full chiral Standard Model, Lorentzian reconstruction,
and a complete nonperturbative BRST Hilbert space remain open
\cite{QFTAudit2026}.

\item[Full AQFT equivalence.] Open.  The corpus has not yet supplied one
selected upper local net together with the naturality interface
\eqref{eq:naturality}, physical state space, covariance, spectrum condition,
time-slice theorem, continuum completion, and target-equivalence certificate.
\end{description}

\section{Scoped reconstruction theorem}

\begin{theorem}[Scoped MTT--AQFT relation]\label{thm:scoped}
The corrected MTT corpus supports the following implications:
\begin{enumerate}
\item admissible contexts can carry a pregeometric precosheaf when their
algebras and functorial extension maps are supplied;
\item a selected Lorentzian base, a local upper net, and a decomposable
coherent projector produce the physical local net
\eqref{eq:compressed-net} on the compatible subalgebra;
\item upper isotony and Einstein causality descend to that physical net; and
\item a natural chart-to-region interface identifies the pregeometric and
physical nets only to the degree expressed by its component maps.
\end{enumerate}
No commutation relation follows from non-joint representability.  No absence
of an abstract quasilocal algebra, state, faithful abstract representation,
horizon law, or temporal arrow follows from absence of a global admissible
chart.
\end{theorem}

\begin{proof}
The first item is the definition and functoriality requirement of
$\mathfrak B$.  The second and third are Theorem~\ref{thm:descent}.  The fourth
is Proposition~\ref{prop:naturality}.  The final exclusions follow from the
typing arguments and the inductive-limit state proposition above.
\end{proof}

\section{Changes from Version 1}

Version~2 makes the following theorem-level corrections.
\begin{enumerate}
\item It replaces the chart-only AQFT derivation by separate pregeometric and
physical constructions.
\item It removes the inference from failure of joint representability to
commutation.
\item It removes automatic observable extension and requires declared
functorial morphisms.
\item It derives physical isotony from upper algebra inclusion and coherent
compression.
\item It derives physical locality from upper spacelike commutation, not from
chart non-comparability.
\item It restores the possible quasilocal inductive limit and distinguishes it
from a single global admissible chart.
\item It replaces claims of automatic horizons and irreversibility by their
actual state, causal, and dynamical theorem requirements.
\item It adds the chart-to-region naturality gate and updates the QFT status to
the current conditional/finite-domain frontier.
\end{enumerate}

\section{Conclusion}

MTT currently has a rigorous but conditional AQFT bridge.  A local upper
theory can be compressed to a coherent physical sector without losing isotony
or spacelike commutation, provided the reduction is fiberwise and the
observables preserve that sector.  This is meaningful progress: it shows that
coherent reduction need not destroy an already established local algebraic
structure.

The result does not make locality emerge from non-comparability.  Nor does it
construct the upper QFT, the physical state, or the chart-to-region map.
Separating those obligations gives the program a precise next target: select
one local Lorentzian upper theory from MTT data and prove the naturality and
equivalence conditions linking its compressed net to the admissibility
precosheaf.

\begin{thebibliography}{9}

\bibitem{HaagKastler1964}
R. Haag and D. Kastler,
\emph{An Algebraic Approach to Quantum Field Theory},
Journal of Mathematical Physics \textbf{5} (1964), 848--861,
doi:10.1063/1.1704187.

\bibitem{BFV2003}
R. Brunetti, K. Fredenhagen, and R. Verch,
\emph{The Generally Covariant Locality Principle---A New Paradigm for Local
Quantum Physics},
Communications in Mathematical Physics \textbf{237} (2003), 31--68,
arXiv:math-ph/0112041.

\bibitem{ChilianFredenhagen2008}
B. Chilian and K. Fredenhagen,
\emph{The Time Slice Axiom in Perturbative Quantum Field Theory on Globally
Hyperbolic Spacetimes},
arXiv:0802.1642.

\bibitem{Projectionv2}
P. Nero,
\emph{The Projection--Admissibility Principle: Descent, Recovery, and
Structural Constraints on Effective Description}, version 2, 2026.

\bibitem{FPVIv4}
P. Nero,
\emph{Fixed Points VI: Formal Synthesis and Physical Interpretations},
corrected version 4, 2026.

\bibitem{Atlasv3}
P. Nero,
\emph{Modal Triplet Theory: A Typed Relationship Atlas}, version 3, 2026.

\bibitem{QFTAudit2026}
P. Nero,
\emph{MTT Selected Quantization and Nonperturbative-QFT Strict-Upgrade Audit},
technical audit version 1, 2026.

\end{thebibliography}

\end{document}
